\section{Truncations, coordinates, and maximal averages}
\label{sec:preliminaries}

We first fix the representatives of all finite truncations and the
coordinate convention for their linearizations. We also record the
one-dimensional maximal estimates that will control changes of the
matrix energy. Constants denoted by $C$ may change from line to line;
their additional dependences are indicated explicitly.
For an integrable function $f$ and $L>0$, set
\[
 \Dil_Lf(t)=L^{-1}f(t/L).
\]
All inputs and matrices may be complex. Scalar forms are complex
multilinear, so pairings with dual inputs do not include a conjugate.

\begin{lemma}[Finite truncations]\label{lem:finite-truncations}
For $F,G\in L^3(\R^2)$, all finite annular integrals defining
$B_{\varepsilon,R}(F,G)$ are absolutely convergent outside one common
null set. On its complement they depend continuously on $(\varepsilon,R)$ in
$0<\varepsilon<R<\infty$. Changing representatives changes this family
only on a common null set. Moreover,
\begin{equation}\label{eq:finite-annulus-bound}
 \left\|\int_{\varepsilon<|t|<R}
 |F(x+t,y)G(x,y+t)|\,\frac{\dd t}{|t|}\right\|_{3/2}
 \le 2\log(R/\varepsilon)\|F\|_3\|G\|_3.
\end{equation}
\end{lemma}
\begin{proof}
For each fixed $t$, H\"older's inequality and translation invariance
bound the $L^{3/2}$ norm of the product by $\|F\|_3\|G\|_3$.
Minkowski's integral inequality gives \eqref{eq:finite-annulus-bound}.
Apply this estimate to the countably many annuli $1/m<|t|<m$, $m\ge2$.
Their union of exceptional null sets works for every finite annulus.
Absolute continuity of the integral gives endpoint continuity.
The same estimate applied to a zero $L^3$ difference proves the assertion
about representatives, first on those countably many annuli and hence
on all finite annuli.
\end{proof}

In particular, suprema over rational endpoint pairs agree with suprema
over all endpoint pairs whenever only finitely many endpoints occur
in a choice. Rational exhaustion makes the variation and the count
quantities in the introduction measurable.

\subsection{The triangular coordinates}

For most of the proof the inputs are smooth and compactly supported.
The change of variables
\[
 x=-a,\qquad y=-b,\qquad t=a+b+c
\]
has absolute determinant one. Define
\begin{equation}\label{eq:pair-maps}
 G_1(b,c)=F(b+c,-b),\qquad G_2(c,a)=G(-a,c+a).
\end{equation}
Both pair maps preserve $L^3$ norms. With a dual input $G_0(a,b)$,
the expressions to be estimated take the form
\begin{equation}\label{eq:triangular-form}
 \Lambda_{\mathcal K}(G_0,G_1,G_2)
 =\iiint G_0(a,b)G_1(b,c)G_2(c,a)
       \mathcal K_{a,b}(a+b+c)\dd a\dd b\dd c.
\end{equation}
The maps in \eqref{eq:pair-maps} are invertible, so estimates for
arbitrary smooth compactly supported $G_v$ are equivalent to the
corresponding dual estimates in the original coordinates.

A \emph{step choice} is a choice from a finite menu that is constant on
each cell of a partition of $\R^2$ by finitely many horizontal and
vertical lines. Unbounded cells are allowed. For annular linearizations,
each menu item is a list of at most $n$ disjoint radius intervals
$(\varepsilon_j,R_j)$, together with coefficients $|\beta_j|\le1$.
At a fixed output point the intervals are disjoint up to endpoints.
We first prove all estimates for step choices. Values on cell boundaries
can be assigned any menu item; they do not affect the integrals.

For such choices and smooth compactly supported inputs, finite-mesh
estimates pass to \eqref{eq:triangular-form} by ordinary Riemann sums.
Indeed, choose $T$ so large that all coordinates in the input supports
lie in $[-T,T]$, put $h=T/N$, and sample at the nodes $u_i=ih$.
After any auxiliary frequency integration over a fixed bounded block,
the finitely many kernels used below are continuous in their scalar
variable. The resulting integrands on $[-T,T]^3$ are bounded and
continuous off finitely many coordinate planes. Their Riemann sums
converge, as do the discrete $L^3$ norm sums. No limit of matrix
energies is taken.

\subsection{Maximal averages on a line}

For functions on $\R$ and sequences on $\mathbb Z$, respectively, let
\[
 M_{\R}f(x)=\sup_{r>0}\frac1{2r}\int_{x-r}^{x+r}|f(t)|\dd t,
 \qquad
 M_{\mathbb Z}f(j)=\sup_{r\in\mathbb Z_{\ge0}}
       \frac1{2r+1}\sum_{|m-j|\le r}|f(m)|.
\]
The discrete operator includes singleton averages. We use these
operators on individual coordinates of a function or array as well.

\begin{lemma}[Line maximal estimates]\label{lem:line-maximal}
On either space, the centered maximal operator satisfies
\[
 \mu\{Mf>\rho\}\le\frac3\rho\|f\|_1,
 \qquad
 \|Mf\|_p^p\le\frac{6p\,2^{p-1}}{p-1}\|f\|_p^p
 \quad(1<p<\infty).
\]
Let $g_\sigma(t)=(2\pi\sigma)^{-1/2}e^{-t^2/(2\sigma)}$.
There is an absolute constant $C$ such that
\begin{align}
 \sum_m h g_\sigma((m-j)h)|f(m)|
   &\le C M_{\mathbb Z}f(j), &&\sigma\ge h^2,
   \label{eq:discrete-gaussian-maximal}\\
 \int g_\sigma(t)|f(x-t)|\dd t
   &\le C M_{\R}f(x), &&\sigma>0.
   \label{eq:continuous-gaussian-maximal}
\end{align}
\end{lemma}
\begin{proof}
We recall the covering proof of the Hardy--Littlewood estimates
\cite{HardyLittlewood1930}. From any finite family of intervals, repeatedly
select one of largest radius and discard those meeting it. The selected
intervals are disjoint, and the union of the family lies in their
triple-radius enlargements. Each enlargement has at most three times
the original measure, also for integer intervals. Apply this to witnessing
intervals centered at a finite subset of a discrete level set, or to a
finite witnessing cover of a compact subset of a continuous level set.
The latter level set is open because fixed-radius averages are continuous.
Taking suprema over the finite or compact subsets proves the weak bound.

For $f\in L^p$, split at height $\rho/2$. The low part has maximal
function at most $\rho/2$, while the high part is integrable. Thus
\[
 \mu\{Mf>\rho\}
 \le\frac6\rho\int_{|f|>\rho/2}|f|\dd\mu.
\]
The distribution formula and Tonelli's theorem give the strong bound
after integrating $p\rho^{p-1}\dd\rho$. Truncation of the distribution
integral first makes this argument independent of any a priori
finiteness assertion.

For \eqref{eq:discrete-gaussian-maximal}, put $a=\sqrt\sigma/h\ge1$.
The kernel is $(\sqrt{2\pi}a)^{-1}e^{-(m-j)^2/(2a^2)}$.
Split the sum into $|m-j|\le a$ and the annuli
$2^{l-1}a<|m-j|\le2^la$, $l\ge1$. The sum of $|f|$ in the
corresponding ball is at most $3\cdot2^la M_{\mathbb Z}f(j)$.
The resulting series is bounded by a constant times
$1+\sum_{l\ge1}2^l e^{-4^{l-1}/2}$.
The same shell argument with integrals proves
\eqref{eq:continuous-gaussian-maximal}.
\end{proof}
